\documentclass[11pt]{article}
\usepackage[a4paper,margin=1in]{geometry}
\usepackage[T1]{fontenc}
\usepackage{lmodern}
\usepackage{microtype}
\usepackage{amsmath,amssymb,amsthm,mathtools}
\usepackage{booktabs}
\usepackage{xcolor}
\usepackage[numbers]{natbib}
\usepackage[colorlinks=true,allcolors=blue]{hyperref}

\newtheorem{theorem}{Theorem}
\newtheorem{proposition}[theorem]{Proposition}

\title{Unbounded Shell-Multiplicity Pathology}
\author{Bin Wang}
\date{July 2026}

\begin{document}
\maketitle

\begin{abstract}
We relax the RH-248 single-use shell zonotope to the full nonnegative cone.
Six of 32 frozen endpoints remain outside the anchored tolerance, with
primal--dual certificates.  The other 26 can be fitted only by shell weights
whose minimum necessary maxima range from $40.58$ to $5.80\times10^{10}$.
Such arbitrary real multiplicities are not a spectral cloud.  The formal
success is therefore a moment-reweighting pathology, not coefficient
identification.
\end{abstract}

\section{Cone relaxation and dual}

Retain the real shell power matrix $V$, target difference $d$, and weights
$\omega_n=1/n$ from RH-248 \citep{WangRH248}.  Remove the upper bounds on
shell use and define
\begin{equation}\label{eq:cone}
 \delta_+=\min_{w\ge0}\sum_{n=2}^{12}omega_n|d_n-(Vw)_n|.
\end{equation}

\begin{theorem}[nonnegative-cone dual]\label{thm:dual}
The cone distance satisfies
\begin{equation}\label{eq:dual}
 \delta_+=
 \max_{\substack{|y_n|\le\omega_n\\V^Ty\le0}}y^Td.
\end{equation}
Thus any feasible dual vector with value greater than $\sigma$ excludes all
nonnegative shell weights, regardless of their size.
\end{theorem}

\begin{proof}
Dualize the weighted $\ell^1$ residual as in RH-248.  The infimum of
$-y^TVw$ over $w\ge0$ is finite exactly when $V^Ty\le0$, in which case it is
zero.  Finite-dimensional LP strong duality gives \eqref{eq:dual}.
\end{proof}

For endpoints with $\delta_+\le\sigma$, a second LP quantifies the pathology:
\begin{equation}\label{eq:cap}
 W_*=\min\{W:\ 0\le w_j\le W,
 \ \sum_n\omega_n|d_n-(Vw)_n|\le\sigma\}.
\end{equation}
This is the smallest possible maximum shell weight, not the weight of an
arbitrary optimal cone solution.

\section{Finite result}

The frozen RH-222 windows and RH-243 anchor give 26 cone passes and six
failures \citep{WangRH222,WangRH243}.  The failing endpoints are
\[
 (0.04,L),\ (0.02,L),\ (0.02,R),\ (0.016,R),\
 (0.00625,L),\ (0.00625,R).
\]
Their smallest distance-minus-tolerance margin is
$2.826242139956541\times10^{-4}$, while the maximum recorded primal--dual
gap is $4.08\times10^{-7}$.  Hence the finite separation is much larger than
the LP consistency error.

For the 26 formal passes, solving \eqref{eq:cap} gives
\[
 40.58443731031147\le W_*
 \le58018432630.629776.
\]
Seventeen cone fits can drive the order-12 jet objective to floating zero,
yet may still require thousands to billions of repetitions.  With the common
cap $W=40$, no endpoint passes; at $W=41$, exactly one endpoint passes.

\begin{center}
\small
\begin{tabular}{lr}
\toprule
quantity & value\\
\midrule
unbounded cone passes/failures & 26 / 6\\
minimum required cap among passes & 40.5844373\\
maximum required cap among passes & $5.8018433\times10^{10}$\\
passes with cap 40 & 0/32\\
passes with cap 41 & 1/32\\
\bottomrule
\end{tabular}
\end{center}

\section{Interpretation and boundary}

An algebraic spectral multiset has fixed nonnegative integer multiplicities.
The variables in \eqref{eq:cone} are arbitrary real repetitions of already
resolved shells, often enormous.  They alter the moment problem without
constructing a corresponding noisy operator or invariant root space.  The 26
LP passes therefore do not supply the RH-243 coefficient bridge.

The six failures and multiplicity explosion concern only the frozen
candidate windows and orders 2--12.  They do not exclude expanded windows,
deeper roots, signed/complex quotient grouping, or continuum mechanisms.  No
all-order envelope follows.  Gate A remains open and Gates B--E are untouched;
no Hilbert--Polya operator, zeta-divisor equality, Riemann-zero
identification, or RH implication is claimed.

\bibliographystyle{plainnat}
\bibliography{references}
\end{document}
